\section{Representatives and classes on the families}

\begin{definition}[Representative and class of a component]
\label{def:c12-class}
Let $\mathfrak F$ be a family of colour-singlet local functionals of degree $0$, stable under $\Lambda$ and under the classical transformations along a subalgebra $\mathfrak h\subset\mathfrak{psl}(4|4)$, with $\hbar L_1^{(0)}=0$ on $\mathfrak F$ for a finite part $\psi$ (Definition \ref{def:c10-hyp} (ii)), and let $H$ be the cohomology of $s$ on $\mathfrak F$. Write $\Lambda^{(1)}(x)=\Lambda^{(1)}_0(x)+\hbar\Lambda^{(1)}_1(x)$ with $\Lambda^{(1)}_0(x)=\pm\{S^{(1)}(x),-\}$ (Proposition \ref{prop:c11-relations}). The representative of the component along $x$ on $O\in\mathfrak F$ is the local functional $\hbar\Lambda^{(1)}_1(x)O=\hbar L_1(O)^{(1)}(x)$ of \eqref{eq:c12-master}; its class is the map $\delta(x)$ induced on $H$ by $-\hbar\Lambda^{(1)}_1(x)$, and $\rho_0(x)$ is the map induced by $-\Lambda^{(1)}_0(x)$. A change of the finite part of the operators or a field redefinition changes $\delta$ by the coboundary $x\mapsto[\rho_0(x),\phi]$ with $\phi\in\mathrm{End}(H)$ preserving the classical Weyl weight, the coupling direction leaves $\delta$ unchanged, and the outer direction $G_B$ changes it by $-a\,\rho_0\circ B$ (Theorem \ref{thm:c12-opalg} (iv)). The scheme-independent class of the components is the class of $\delta$ in $H^1(\mathfrak h;\mathrm{End}(H))$ modulo $\mathbb C\,\rho_0\circ B$; it contains the spectrum of $\delta(\Delta)$ on each classical weight space of $H$.
\end{definition}

\begin{theorem}[The components on the scalar bilinears]
\label{thm:c12-F0}
Let $O=\int\omega\,(\iota_{w_1}A,\iota_{w_2}A)$ with $w_1,w_2\in W_{\mathbb R}^\perp\otimes\mathbb C$, $\psi=0$, $\xi\in\Sigma$ of either chirality, and $\mathcal B_1,\dots,\mathcal B_4$ the functionals of Lemma \ref{lem:c11-cocycle}: $\mathcal B_1=\{S^{(1)}(\xi),O\}$, $\mathcal B_2=(w_1,w_2)\int\omega\,(\rho(A|_{W_{\mathbb R}^\perp})\xi,\lambda)$, $\mathcal B_3=(w_1,w_2)\int\omega\,(\rho(A|_V)\xi,\lambda)$, $\mathcal B_4=(w_1,w_2)\int\omega\,(c,(\xi,\lambda^*))$. With the uniform coefficients of Theorem \ref{thm:c12-ar2}, the representatives are
\begin{equation}
\begin{aligned}
Q_1(O)&=\Big[\big(\tfrac{79}{12}-3\log\tfrac32-t_\xi\big)\mathcal B_1+4\log\tfrac32\,\mathcal B_2+4\log\tfrac32\,\mathcal B_3+0\cdot\mathcal B_4\Big]\ \text{units},\qquad \hbar L_1(O)^{(0)}=-\big(2+4\log\tfrac32\big)(w_1,w_2)\int\omega\,(dc,A|_V)\ \text{units},\\
S_1(O)&=\big(\tfrac{79}{12}-3\log\tfrac32-t_\eta\big)\{S^{(1)}(\eta),O\}\ \text{units}+x_O^\mu\Big[4\log\tfrac32\,\mathcal B_2+4\log\tfrac32\,\mathcal B_3\Big](\rho(e_\mu)\eta)\ \text{units},\\
D_1(O)&=\big(-4t_\Delta-\tfrac43\big)O\ \text{units (traceless)},\qquad \big(-4t_\Delta-\tfrac{148}3\big)O\ \text{units (trace)},\qquad\text{at the point of the density},\\
P_1(O)&=-\big(\tfrac13+t_P\big)\,\partial_bO\ \text{units},\qquad M_1(O)=-\big(\tfrac13+t_M\big)\{S^{(1)}(M_\omega),O\}\ \text{units},\qquad R_1(O)=0,\\
K_1(O)&=v(x_O)^\mu P_1(O)\big|_{P_{e_\mu}}-(b,x_O)\,D^0_1(O)+M^0_1(O)\big|_{\omega_{x_O}}+K^0_1(O),
\end{aligned}
\label{eq:c12-F0}
\end{equation}
with the two-body part of $K^0_1$ equal to $0$ and its one-body part
\[
K_1(O)\big|_{\text{one-body}}=\tfrac{44}3\int\omega\,(b,x)(\iota_{w_1}A,\iota_{w_2}A)\ \text{units}-\theta^K_\perp\big\{S^{(1)}(K_b)_\perp,O\big\}\ \text{units}-\sum_{j=1}^4\kappa^K_j\{G^K_j,O\},
\]
whose part linear in $x$ is $-(b,x_O)$ times the one-body part $-\tfrac{44}3O$ units of $D^0_1$. The coefficient of $\mathcal B_1$ in $Q_1$ is the two-body $-2\log\tfrac32$ plus the one-body $-\big(\log\tfrac32-\tfrac{79}{12}+t_\xi\big)$; the intrinsic part of $S_1$ vanishes. At $t_\Delta=t_M=-\tfrac13$: $D_1=0$ on the traceless bilinears, $D_1=-48\,O$ units $=\tfrac{3g^2N}{2\pi^2}\,O$ on the trace, and $M_1=0$.
\end{theorem}

\begin{proof}
Theorem \ref{thm:c12-master} with the two-body kernels of the pair $(A|_{W_{\mathbb R}^\perp},A|_{W_{\mathbb R}^\perp})$ and the one-body terms from Theorem \ref{thm:c12-sol}; on $O$ only the terms of $\hbar F_1^{(1)}(x)$ with a leg $A^*|_{W_{\mathbb R}^\perp}$ act. Along $\xi$: of $\hbar F_1^{(1)}(\xi)|_2=(\log\tfrac32-\tfrac{79}{12})(T_1+T_2)+t_\xi(S^V_2+S^\perp_2)(\xi)$ units the terms with that leg are $T_2$ in both, and $\{T_2,O\}=\{S^{(1)}(\xi),O\}=\mathcal B_1$, the arrow $A|_{W_{\mathbb R}^\perp}\to\lambda$ of $S^{(1)}(\xi)$ being $T_2$ alone; the two-body part is \eqref{eq:c11-Q1O} (leg sets with $D=0$, Remark \ref{rem:c12-finitesp}) and $L_1^{(0)}(O)$ is Theorem \ref{thm:c11-bilinears} (iii). Along $\eta$ the position part of the solution is $x^\mu\hbar F_1^{(1)}(\rho(e_\mu)\eta)$ with the uniform $t_\eta$ (Theorem \ref{thm:c12-ar2}), $s_0$ and $\int(\rho(A|_V)\eta,\lambda^*)$ have no leg $A^*$, and Corollary \ref{cor:c12-SK} gives the two-body part. Along $\Delta$, $-\{\tfrac{22}3\int(A,A^*),O\}=-\tfrac{44}3O$ and $-\{t_\Delta S^{(1)}(\Delta)_\perp,O\}=-4t_\Delta O$ units at the point, and the two-body sums are $\tfrac{40}3O$ and $-\tfrac{104}3O$ units (Remark \ref{rem:c12-pair}). Along $P_b$, $\Upsilon_{P_b}=-\tfrac13\partial_bO$ units and $-\{t_PS^{(1)}(P_b)_\perp,O\}=-t_P\partial_bO$ units. Along $M_\omega$, $\Upsilon^0_{M_\omega}=0$ on the pair, the position part is $(\omega x_O)^\mu\Upsilon_{P_{e_\mu}}$ (Theorem \ref{thm:c12-pos}), the spin term of $\hbar F_1^{(1)}(M_\omega)$ has no leg $A^*|_{W_{\mathbb R}^\perp}$, and $\{S^{(1)}(M_\omega),O\}$ is the transport $(\omega x)\cdot\partial O$. Along $R_r$, $\Upsilon_{R_r}=0$ and $t_R=0$. Along $K_b$, Corollary \ref{cor:c12-SK} with the $x$-linear part of $\hbar F_1^{(1)}(K_b)|_2$ (Theorem \ref{thm:c12-sol}). The values at $t_\Delta=-\tfrac13$: $-\tfrac43+\tfrac43=0$ and $-\tfrac{148}3+\tfrac43=-48$, with the unit $-\tfrac{g^2N}{32\pi^2}$.
\end{proof}

\begin{corollary}[Classes on the scalar bilinears]
\label{cor:c12-F0inv}
\begin{itemize}
\item (i) The spectrum of $\delta(\Delta)$ on the scalar bilinears is $0$ on the twenty traceless bilinears and $-48$ units $=\tfrac{3g^2N}{2\pi^2}$ on the trace, the Konishi value $2\gamma_K$ with $\gamma_K=\tfrac{3g^2N}{4\pi^2}$: it is independent of the scheme, $t_\Delta=-\tfrac13$ being fixed by the arity-two equations (Theorem \ref{thm:c12-ar2}).
\item (ii) On the traceless bilinears $\delta(\xi)$ is $c_Q\,\rho_0(\xi)$ with $c_Q=\tfrac{79}{12}-3\log\tfrac32-t_\xi$ units, $\delta(\eta)$ is $c_S\,\rho_0(\eta)$ with $c_S=\tfrac{79}{12}-3\log\tfrac32-t_\eta$ units, and $\delta(P_b)=p_1\rho_0(P_b)$ with $p_1=-(\tfrac13+t_P)$ units. The outer direction shifts $c_Q(\xi_\pm)$ by $\mp a$, and a finite part of the operators shifts $c_Q$, $c_S$, $p_1$ by coboundaries; beyond $\delta(\Delta)$, the scheme-independent content of these components on the scalar bilinears is contained in the relations (R2), (R3) of Corollary \ref{cor:c12-rel}.
\item (iii) At $\psi=0$, as functions of the action-side free directions,
\begin{equation}
c_Q(\xi_+)+c_Q(\xi_-)-p_1=c_Q(\xi_+)+c_S(\eta_+)-\tfrac14D_1\big|_{\text{traceless}}=\tfrac{27}4-8\log\tfrac32\ \text{units},
\label{eq:c12-F0inv}
\end{equation}
independent of $t_\Delta$, $t_P$, $t_\xi$, $t_\eta$ and of the inner and outer directions, by \eqref{eq:c12-t}. A finite part of the operators shifts $c_Q$, $c_S$ and $p_1$ by coboundaries, so these combinations depend on the operator scheme; the relations (R2), (R3) of Corollary \ref{cor:c12-rel} on the scalar bilinears contain them together with the components on the descendants of Weyl weight $\tfrac52$ and the arity-two kernels.
\item (iv) The coefficients $c_Q$, $c_S$, every coefficient with $\log\tfrac32$, $P_1$, $L_1^{(0)}(O)$ (removed by $\psi(O)=-(1+2\log\tfrac32)(w_1,w_2)\int\omega\,(A|_V,A|_V)$ units), $M_1$, $R_1$ and the individual pair kernels depend on the scheme.
\end{itemize}
\end{corollary}

\begin{proof}
(i) Theorem \ref{thm:c12-F0} and Theorem \ref{thm:c8-invariants}. (ii) Theorem \ref{thm:c12-F0}, Theorem \ref{thm:c12-opalg} (iv) and Proposition \ref{prop:c11-freedom} (ii), (iv). (iii) $2\cdot\tfrac{79}{12}-6\log\tfrac32-(t_{\xi_+}+t_{\xi_-})+\tfrac13+t_P=\tfrac{81}{12}-8\log\tfrac32$ by \eqref{eq:c12-t}, and $2\cdot\tfrac{79}{12}-6\log\tfrac32-(t_{\xi_+}+t_{\eta_+})+\tfrac13+t_\Delta$ is the same. (iv) Theorem \ref{thm:c11-bilinears} (iii), Corollary \ref{cor:c11-theta} and Theorem \ref{thm:c12-opalg} (iv).
\end{proof}

\begin{proposition}[The twisted differential on the chiral bilinears]
\label{prop:c12-twist}
Let $\xi_\alpha{}^i=u_\alpha z^i$ be of rank one and $w_1,w_2$ with $\rho(w_1)\xi=\rho(w_2)\xi=0$, so that $O=\int\omega\,(\iota_{w_1}A,\iota_{w_2}A)$ is a combination of the bilinears $\int\omega\,(A_{ij}z^j,A_{kl}z^l)$ (Proposition \ref{prop:c11-plane}). Then the order-$\hbar$ part of the twisted differential vanishes on $O$: $\hbar Q^\xi_1(O)=\hbar L_1(O)^{(0)}+Q_1(O)-\tfrac12\hbar L_1(O)^{(2)}(\xi,\xi)=0$ (Definition \ref{def:c8-twisted}), for every value of the action-side free directions, of $\psi$ and of the outer and inner directions. For every quasi-isomorphism $\Phi$ of Theorem \ref{thm:c11-twistcheck}, whatever its normalisation, the class $\Phi\circ[Q^\xi_1]\circ\Phi^{-1}$ vanishes on the image of $[O]$, the classes of the two-letter chiral operators $\operatorname{tr}\gamma^{(I}\gamma^{J)}$.
\end{proposition}

\begin{proof}
$\mathcal B_1=\{S^{(1)}(\xi),O\}$ contains $\iota_{w_a}\Gamma(\xi,\lambda)=(\rho(w_a)\xi,\lambda)=0$ [ESW, (28)], and $(w_1,w_2)=0$ (Proposition \ref{prop:c11-plane}) removes $\mathcal B_2$, $\mathcal B_3$, $\mathcal B_4$ and $L_1^{(0)}(O)$; Theorem \ref{thm:c12-F0} holds for every $\xi$, being linear in $\xi$, and its one-body term is a multiple of $\mathcal B_1$ for every value of $t_\xi$. The $(\xi,\xi)$ component of a functional of degree $0$ and weight $2$ has degree $-1$ and weight $3$, so a two-letter output contains an antifield of weight at least $\tfrac52$ and a field of weight at least $1$, total at least $\tfrac72>3$; the one-body term $-\{\hbar F_1^{(2)}(\xi,\xi),O\}$ vanishes, $\hbar F_1^{(2)}(\xi,\xi)$ having the legs $(\lambda^*,\lambda^*)$ (Table \ref{tab:c12-legs} and Theorem \ref{thm:c12-bound}). Under the scheme changes $Q^\xi_1$ changes by a commutator with $Q^\xi_0$ or by a multiple of $Q^\xi_0$ (Theorem \ref{thm:c8-twist}), both zero on the class of $O$, and $\Phi$ is linear.
\end{proof}

\begin{remark}[The comparison with the holomorphic twist]
\label{rem:c12-twistcmp}
Proposition \ref{prop:c12-twist} is the case $n=2$ of [BK, Sec. 3.4.1, (3.77)], $Q_1(\gamma^I\gamma^J)=0$, at the level of classes. The full comparison of $[Q^\xi_1]$ modulo $\operatorname{Im}Q^\xi_0+\mathbb C\,S^\xi$ with [BK, (3.7), (3.59)--(3.70)] is the statement of Theorem \ref{thm:c11-twistcheck}; its inputs are the transport map $\Phi$ of [ESW, Thm 10.9] with the normalisation of [BK] for the letters $b$, $\beta$, $\gamma$, $\partial_{\dot\alpha}c$, the arity-one entries of Table \ref{tab:c11-twistentries}, which are kernels of the form \eqref{eq:c12-kernel} (their leg sets have $D=0$), the arity-two kernels along $(\xi,\xi)$ for the entries $b\,c$ and $b\,b$ with $b$ represented by the field strength, and the finite parts $\psi$, $t_\xi$, $a$ that move the colour-order-dependent terms in $\log\tfrac32$ of the individual entries into $\operatorname{Im}Q^\xi_0$. The comparison is a cross-check: no number of the components is taken from it.
\end{remark}

\begin{proposition}[The dilatation component on the families]
\label{prop:c12-spectra}
On the families of Definition \ref{def:c11-families}, at $t_\Delta=-\tfrac13$ and $\psi$ as in Definition \ref{def:c10-hyp} (ii), $\delta(\Delta)$ has the spectra of Table \ref{tab:c12-spectra}, independent of the scheme; with the unit $-\tfrac{g^2N}{32\pi^2}$, $D_1=-8$ units times the eigenvalue $E$ of $\sum_lH_{l,l+1}$ is $\tfrac{g^2N}{4\pi^2}E=2\gamma$. On twist two, the fermion $\mathfrak{sl}(2)$ sector and the three families with gaugino or field-strength letters, $\sum_{a<b}\pm\Upsilon_{P_n}=-\tfrac13\partial_n$ units, so $P_1=-(\tfrac13+t_P)\partial_n$ units; on twist two, the fermion sector, the family with two gaugino components and the mixed family the dilatation commutes with the classical translation along the null direction, $D_1^{(S'+1)}\partial_n=\partial_nD_1^{(S')}$, as an exact matrix identity. On every complete row of the single-pair tables of Remark \ref{rem:c12-method} the pair form of Remark \ref{rem:c12-pair} holds.
\end{proposition}

\begin{tabular}{l|l|l}
family & range & spectrum of $D_1$ in units\\
\hline
$\mathfrak{su}(2)$ sector $\operatorname{tr}(Z^{L-M}X^M)$ & $L\le8$ & $-8\sum_l(1-P_{l,l+1})$, the Heisenberg spectrum\\
$\mathfrak{su}(1|1)$ sector $\operatorname{tr}(Z^{L-M}\psi^M)$ & $L\le6$ & $-8\sum_l(1-\Pi_{l,l+1})$, graded exchange $\Pi$; Table \ref{tab:c11-su11} times $-8$\\
twist two $\operatorname{tr}(D_n^kZ\,D_n^{S'-k}Z)$ & $S'\le6$ & $-32h(j)$, $j\le S'$ even: $0,-48,-\tfrac{200}3,-\tfrac{392}5$\\
fermion $\operatorname{tr}(D_n^k\psi\,D_n^{S'-k}\psi)$ & $S'\le5$ & $-32h(j+1)$, $j\le S'$ odd\\
two gaugino components with $Z'F_{11}$ & $S'\le3$ & $0,-48$; $0,(-48)^3$; $0,(-48)^4,-\tfrac{200}3$; $0,(-48)^4,(-\tfrac{200}3)^3$\\
gluon $\operatorname{tr}(D_n^kF_{11}\,D_n^{S'-k}F_{11})$ & $S'\le2$ & $-48$; $-48$; $-\tfrac{200}3,-48$\\
mixed $\operatorname{tr}(D_n^kF_{11}\,D_n^{S'-k}\psi)$ & $S'\le4$ & $-48$; $(-48)^2$; $-\tfrac{200}3,(-48)^2$; $(-\tfrac{200}3)^2,(-48)^2$; $-\tfrac{392}5,(-\tfrac{200}3)^2,(-48)^2$\\
\end{tabular}
\tabcaption{The spectra of $D_1$ on the families, $h$ the harmonic number and a power the multiplicity; the entries separated by semicolons are $S'=0,1,\dots$. The first four rows are computed by both evaluations of Remark \ref{rem:c12-method}, the last three by the first, with the eigenvalues recomputed independently.\label{tab:c12-spectra}}

\begin{proof}
Theorem \ref{thm:c12-D1} with the pair kernels of the families, evaluated through the two-field parts (Proposition \ref{prop:c10-twist} and the remark before Proposition \ref{prop:c11-targets}); the one-body term at $t_\Delta=-\tfrac13$ and the composite pair kernels give the spectra (Theorem \ref{thm:c11-families}). Scheme independence is Theorem \ref{thm:c8-invariants} (iii), $t_\Delta$ being fixed by Theorem \ref{thm:c12-ar2}. The translation identity is (R7) of Corollary \ref{cor:c12-rel} along $P_n$.
\end{proof}

\begin{remark}[The spectra as a cross-check]
\label{rem:c12-specref}
The spectra of Table \ref{tab:c12-spectra} are $-8$ units times the eigenvalues of Proposition \ref{prop:c11-targets}, of Proposition \ref{prop:c10-su2} and of Table \ref{tab:c11-su11}, the one-loop dilatation operator of [B03c]: $\gamma=\tfrac{g^2N}{2\pi^2}h(S)$ on twist two, $\gamma=\tfrac{3g^2N}{4\pi^2}$ on $\operatorname{tr}(F_{11}F_{11})$ and the trace bilinear, the Konishi multiplet. The agreement is a cross-check of the construction; the numbers come from \eqref{eq:c12-master}.
\end{remark}

\section{The one-loop algebra on operators}

\begin{proposition}[The arity-two component from the defect flow]
\label{prop:c12-master2}
Let $\hbar F_1$ solve \eqref{eq:c12-eqs} through arity two along $(x,y)$, with the local representatives through arity one along $x$, $y$, $[x,y]$. For $O$ of degree $0$, with $T$ finite if $x$ or $y$ is $K_b$ or if $(x,y)=(\eta_1,\eta_2)$ with $\Gamma(\eta_1,\eta_2)|_V\neq0$, and $T=\infty$ otherwise,
\[
\hbar L_1(O)^{(2)}(x,y)=\sum_{a<b}\pm\Upsilon_{x,y}(X_a,X_b)-\big\{\hbar F_1^{(2)}(x,y),O\big\}+\big(\text{trees of }\hbar F_1^{(1)}(x),\hbar F_1^{(1)}(y)\text{ with }S^{(1)}(y),S^{(1)}(x)\big),
\]
with $\Upsilon_{x,y}$ the coefficient of $\varepsilon^0$ of the cycles through the pair with one special edge and either two vertices $S^{(1)}(x)$, $S^{(1)}(y)$ or one vertex $S^{(2)}(x,y)$; the bridges cancel their counterterm partners as in Lemma \ref{lem:c12-bridge}.
\end{proposition}

\begin{proof}
The proof of Theorem \ref{thm:c12-master} applies in arity two: the defect vertex has the arity-two parts $\pm\{S^{(2)}(x,y),O\}$, the special-edge parts of the products of two arity-one brackets, $-\{\hbar F_1^{(2)}(x,y),O\}_\varepsilon$, $-\{I^{CT,(2)}(x,y),O\}_\varepsilon$ and the arity-two counterterm of $O$; the classes are the same with two arity-one vertices or one arity-two vertex, and (Ac) vanishes by the arity-two equation. The trees with $\hbar F_1^{(1)}(x)$ and $S^{(1)}(y)$ attached to $O$ are the arity-two parts of the products $\Lambda\Lambda$ in the local representative (Theorem \ref{thm:c8-L} (iv)).
\end{proof}

\begin{theorem}[The one-loop algebra]
\label{thm:c12-opalg}
Let $\mathfrak F$, $H$, $\rho_0$, $\delta$ be as in Definition \ref{def:c12-class}.
\begin{itemize}
\item (i) If the quantisation holds through arity one along $x$ and the local representatives exist, then $s\big(\hbar\Lambda^{(1)}_1(x)O\big)=0$ for every $O\in\mathfrak F$ with $sO=0$ (R1).
\item (ii) If moreover the quantisation holds through arity two along $(x,y)$ and the local representatives exist through arity two, then on the cocycles of $\mathfrak F$
\begin{equation}
\big[\Lambda^{(1)}_1(x),\Lambda^{(1)}_0(y)\big]O+\big[\Lambda^{(1)}_0(x),\Lambda^{(1)}_1(y)\big]O+\Lambda^{(1)}_1([x,y])O=\Lambda^{(0)}\big(\chi_{x,y}(O)\big),
\label{eq:c12-Rxy}
\end{equation}
with $\chi_{x,y}(O)$ local of degree $-1$, built from $\hbar L_1(O)^{(2)}(x,y)$ and $\{S^{(2)}(x,y),O\}$ (Corollary \ref{cor:c8-homotopy} (ii)).
\item (iii) Under (ii) for all pairs in $\mathfrak h$, $\delta([x,y])=[\delta(x),\rho_0(y)]+[\rho_0(x),\delta(y)]$, graded commutators: $\delta$ is a one-cocycle of $\mathfrak h$ with values in $\mathrm{End}(H)$.
\item (iv) A finite part $\psi'$ or a field redefinition changes $\delta$ by the coboundary $x\mapsto[\rho_0(x),\psi'_H]$; the coupling direction leaves $\delta$ unchanged; $G_B$ changes it by $-a\,\rho_0\circ B$, a cocycle that on $\mathfrak{psl}(4|4)$ is not a coboundary.
\end{itemize}
\end{theorem}

\begin{proof}
(i), (ii) Theorem \ref{thm:c8-L} (iv), $L\circ L=0$, at arity one and two on an arity-zero input at order $\hbar$, with Corollary \ref{cor:c8-homotopy} and Proposition \ref{prop:c11-relations}; the right side of (ii) is $[\Lambda^{(0)},H_{x,y}]O=\Lambda^{(0)}(H_{x,y}O)$ for $\Lambda^{(0)}O=0$. (iii) On cohomology $\Lambda^{(0)}$-exact terms vanish; $\rho_0$ is a representation by the order-$\hbar^0$ identity, the term $S^{(2)}$ acting as $0$ on cohomology, and the order-$\hbar$ part of $[\rho(x),\rho(y)]=\rho([x,y])$ with $\rho=\rho_0+\delta$ is the cocycle identity. (iv) Proposition \ref{prop:c8-laws}: conjugation by $1+\hbar\psi'$ changes $\rho$ by $[\rho_0,\psi'_H]$ at order $\hbar$; Theorem \ref{thm:c8-invariants} (ii) for $G_B$, and $B$ is outer (Proposition \ref{prop:c8-B}).
\end{proof}

\begin{corollary}[The relations on cohomology]
\label{cor:c12-rel}
Under the hypotheses of Theorem \ref{thm:c12-opalg}, on cohomology:
\begin{itemize}
\item (R2) for $\xi,\xi'\in\Sigma$, $[\delta(\xi),\rho_0(\xi')]+[\rho_0(\xi),\delta(\xi')]=\delta(P_{\Gamma(\xi,\xi')|_V})$, with right side $0$ for $\xi,\xi'\in S_+\otimes W$;
\item (R3) for $\xi\in\Sigma$, $\eta\in\Sigma^*$, $[\delta(\xi),\rho_0(\eta)]+[\rho_0(\xi),\delta(\eta)]=\tfrac12(\xi,\eta)\delta(\Delta)+\delta(M_{\xi\eta})+\delta(R_{\xi\eta})$ (Proposition \ref{prop:c2-special});
\item (R4) for a rank-one $\xi$, $\{\delta(\xi),\rho_0(\xi)\}=0$, the cohomological form of $Q^\xi_0Q^\xi_1+Q^\xi_1Q^\xi_0=0$ (Theorem \ref{thm:c8-twist});
\item (R5) $[\delta(K_b),\rho_0(P_{b'})]+[\rho_0(K_b),\delta(P_{b'})]=(b,b')\delta(\Delta)+\delta(M_{bb'})$;
\item (R6) for $\eta_1,\eta_2\in\Sigma^*$, $[\delta(\eta_1),\rho_0(\eta_2)]+[\rho_0(\eta_1),\delta(\eta_2)]=\delta([\eta_1,\eta_2])$, a special conformal field for $\Gamma(\eta_1,\eta_2)|_V\neq0$;
\item (R7) on a family of operators at one point, graded by the classical Weyl weight and closed under the classical action of the generators considered, $[\delta(\Delta),\rho_0(x)]=0$: the one-loop dilatation commutes with the classical action and is constant on the classical multiplets of the family;
\item (R8) if the family is finite-dimensional and stable under the classical rotations, the restriction of $\delta$ to $\mathfrak{so}(V)\oplus\mathfrak{so}(W_{\mathbb R}^\perp)$ is a coboundary, and a finite part on the family makes $M_1\cong0$ and $R_1\cong0$.
\end{itemize}
(R7) along the null translation $P_n$ holds as an exact matrix identity on twist two, the fermion $\mathfrak{sl}(2)$ sector, the family with two gaugino components and the mixed family, and on the $\mathfrak{su}(2)$ sector.
\end{corollary}

\begin{proof}
(R2)--(R6): Theorem \ref{thm:c12-opalg} (iii) with the brackets of Proposition \ref{prop:c2-special} and Remark \ref{rem:c7-KDelta}. (R7): $[\Delta,x]=\lambda_xx$, so $\lambda_x\delta(x)=[\delta(\Delta),\rho_0(x)]+[\rho_0(\Delta),\delta(x)]$; at a point $\rho_0(\Delta)$ acts on the classical Weyl weight $w$ as a fixed multiple of $w$, and $\delta(x)$ maps weight $w$ to $w-\tfrac12\lambda_x$, as $\rho_0(x)$ does (Proposition \ref{prop:c10-bound1}), so $[\rho_0(\Delta),\delta(x)]=\lambda_x\delta(x)$ and the remaining term vanishes. (R8): $\mathfrak{so}(4,\mathbb C)\oplus\mathfrak{so}(6,\mathbb C)$ is semisimple and $\mathrm{End}(H)$ finite-dimensional, so the first cohomology vanishes (Whitehead's first lemma), and a coboundary $x\mapsto[\rho_0(x),\phi]$ is removed by the finite part $-\phi$. The last statement: Proposition \ref{prop:c12-spectra}.
\end{proof}

\begin{proposition}[Free directions fixed by the algebra on operators]
\label{prop:c12-tfix}
Suppose that on the cohomology of a family the free directions of $\hbar F_1^{(1)}(x)$ act as $-t_x$ units times $\rho_0(x)$, with $t_x$ constant on the blocks $P$, $K$, $\Delta$, $\mathfrak{so}(V)$, $\mathfrak{so}(W_{\mathbb R}^\perp)$, $S_+\otimes W$, $S_-\otimes W^*$, $S_+\otimes W^*$, $S_-\otimes W$; along $\Sigma$, $\Sigma^*$, $P_b$, $\Delta$, $\mathfrak{so}(V)$, $\mathfrak{so}(W_{\mathbb R}^\perp)$ this holds on every family, the sector coefficients being uniform (Theorem \ref{thm:c12-ar2}). If (R2)--(R7) hold for two choices of these numbers, their difference $t$ satisfies $t_{[x,y]}=t_x+t_y$ whenever $\rho_0([x,y])\neq0$ on the cohomology, and $t=aB+\tfrac c2\lambda$: the outer direction and the finite part $\tfrac c2\{S^{(1)}(\Delta),-\}$ (Proposition \ref{prop:c9-scheme}). In particular the difference vanishes on $\Delta$, $\mathfrak{so}(V)$ and $\mathfrak{so}(W_{\mathbb R}^\perp)$.
\end{proposition}

\begin{proof}
With $\delta_t(x)=-t_x\rho_0(x)$ the cocycle identity reads $-t_{[x,y]}\rho_0([x,y])=-(t_x+t_y)\rho_0([x,y])$. From $[\Delta,x]=\lambda_xx$ with $\lambda_x\neq0$, $t_\Delta=0$. Write $\alpha$, $\beta$ for $t$ on $S_+\otimes W$ and $S_-\otimes W^*$; the brackets $[\Sigma,\Sigma^*]\ni\Delta$ of equal chirality give $t=-\alpha$, $-\beta$ on $S_+\otimes W^*$, $S_-\otimes W$ and $t=0$ on the rotations they produce; $[S_+\otimes W,S_-\otimes W^*]\ni P$ gives $t_P=\alpha+\beta$; $[P_b,\eta]=\rho(b)\eta$ and $[K_b,\xi]=\rho(b)\xi$ are consistent with $t_K=-(\alpha+\beta)$ and $[K,P]\ni\Delta$. The two parameters are $a=\tfrac12(\alpha-\beta)$ and $c=-(\alpha+\beta)$, with $B=\pm1$ (Definition \ref{def:c8-B}) and $\lambda=-1,+1,-2,+2$ on $\Sigma,\Sigma^*,P,K$.
\end{proof}

\begin{lemma}[Ward identity of the two-point cochain from the defect flow]
\label{lem:c12-ward2}
Let $O_1$, $O_2$ be gauge-invariant local functionals of degree $0$ with disjoint supports and local representatives with $L(O_a)^{(0)}=0$ modulo $\hbar^2$, and let $\langle\,,\rangle_T$ be the arity-zero two-point function at scale $T$ (Definition \ref{def:c8-2pt}). In arity one along a generator $Y$, modulo $O(\hbar^2)$,
\begin{equation}
\big\langle\Lambda^{(1)}(Y)O_1,O_2\big\rangle_T+(-1)^{|Y||O_1|}\big\langle O_1,\Lambda^{(1)}(Y)O_2\big\rangle_T=B_T(Y;O_1,O_2),
\label{eq:c12-ward2}
\end{equation}
where $B_T$ is $\hbar\Delta_T$ of the quadratic part at the zero field of the arity-one two-point cochain. $B_\infty=0$, and along $\Sigma$, $\Sigma^*$, $\Delta$, the translations and the rotations \eqref{eq:c12-ward2} holds at $T=\infty$ with zero right side.
\end{lemma}

\begin{proof}
By Lemmas \ref{lem:c8-transfer} and \ref{lem:c8-flow} applied to $J+t_1O_1(\varepsilon)+t_2O_2(\varepsilon)$, the $t_1t_2$-part of $d_T$ of the effective functional at scale $T$ is the flow of the special-edge part of $\{t_1O_1(\varepsilon),t_2O_2(\varepsilon)\}_\varepsilon$, of order $e^{-|x_1-x_2|^2/4\varepsilon}$ at separated supports, together with the flows with $t_aL(O_a)$ in place of $t_aO_a$. At the zero field $Q$ and $d_{CE}$ give $0$, $\{F,G\}_T$ needs one-leg parts of both arguments, which vanish for colour singlets and for the effective interaction, and $\hbar\Delta_T$ of the $t_1t_2$-part leaves $B_T$; $\Delta_\infty=0$ (Definition \ref{def:c3-delta}).
\end{proof}

\begin{remark}[$K=P^\dagger$ as a cross-check]
\label{rem:c12-KP}
Assume (H1) a reflection-compatible one-loop quantisation through arity one along $P_n$, $K_n$ and the odd generators considered and through arity two along $(P_n,K_n)$ (Theorem \ref{thm:c8-fields} (iii)); (H2) local representatives on a family $\mathfrak F$ of gauge-invariant cocycles at the origin of one length, stable under $\Lambda^{(1)}(P_n)$, $\Lambda^{(1)}(K_n)$ and the classical generators, with $L_1^{(0)}=0$ for a reflection-compatible finite part; (H3) $\lim_{T\to\infty}B_T(K_n;O_1,O_2)=0$ on the descendants of $\mathfrak F$ entering the radial inner product. Lemma \ref{lem:c12-ward2} gives the Ward identity of Theorem \ref{thm:c8-adjoint} along every generator for which the two-point cochain converges, and along $K_n$ through (H3), $K_1$ being independent of $T$ (Proposition \ref{prop:c12-KT}). Theorem \ref{thm:c8-adjoint} and Corollary \ref{cor:c8-gram} with $P_b^\dagger=-K_{\bar b}$ (Theorem \ref{thm:c8-dag} (i)) then give, on $\mathfrak F$ with the radial inner product of Definition \ref{def:c8-inner},
\[
K_1\ \text{along}\ K_{\bar b}=-\big(P_1\ \text{along}\ P_b\big)^{\rm adj}+\hbar\big[\{S^{(1)}(K_{\bar b}),-\},\mathrm{Gram}_{\rm tree}^{-1}\mathrm{Gram}_1\big],
\]
and the same form for $S_1$ against $Q_1$ and for $D_1$; in a scheme with $P_1=0$ on $\mathfrak F$ (Definition \ref{def:c8-radial}) the first term is absent. In the simplest graph of $B_T(K_n)$ the antifield leg $A^*$ of $S^{(1)}(K_n)$ at $y$, joined by $P\star v.$ to $O_1$, is contracted by $K_T$ with a letter of $O_2$, an integrand of order $|y|^{-1}$, and $\int_{|y|^2\lesssim T}|y|^{-1}T^{-2}d^4y=O(T^{-1/2})$. The identity is a cross-check of $K_1$, which \eqref{eq:c12-master} computes from its own defect cycles at a finite endpoint.
\end{remark}

\section{Former shortcuts, method and continuation}

\begin{remark}[The former shortcuts]
\label{rem:c12-shortcuts}
\begin{itemize}
\item The coefficient $\theta=-\tfrac13$ is computed from the construction twice, from $(\Delta,P_b)$ and from $(\Delta,\xi_\pm)$ (Theorem \ref{thm:c12-ar2}); Theorem \ref{thm:c9-kappa} and Corollary \ref{cor:c11-theta} are cross-checks (Remark \ref{rem:c12-thetacheck}).
\item ``$D_1$ is four times the logarithmic coefficient'' is Theorem \ref{thm:c12-D1}, with Theorem \ref{thm:c12-ward} (i) and Lemma \ref{lem:c12-obeven}; on the action side, $\hbar\,\mathrm{Ob}^{(1)}(\Delta)=4L^{(0)}+$ (composite part) is Theorem \ref{thm:c12-second} (v). On two legs it holds placement class by placement class: the defect cycles at the special edge give $[-4,16,-4]$ units and the line placement $[\tfrac13,-\tfrac13,-\tfrac23]$ units on $\int(dA,dA)$, $\int(\lambda,{\not d}\lambda)$, $\int(dc,A^*)$, whose sum is \eqref{eq:c12-ob1} along $\Delta$.
\item The covariance of the regulator along $P_b$, $M_\omega$, $R_r$ is Proposition \ref{prop:c12-composite} (i) and Theorem \ref{thm:c12-ward} (ii); $\mathrm{Ob}^{(1)}(P_b)|_2=0$ and $\mathrm{Ob}^{(1)}(R_r)=0$ are also computed by direct cycles (Theorem \ref{thm:c12-ob1}).
\item $S_1$ from $Q_1$ by the position decomposition is Theorem \ref{thm:c12-pos} and Corollary \ref{cor:c12-SK}; $S_1$ is computed from its own defect cycles along $\eta$ and from the action-side solution along $\eta$, whose position part is $x^\mu\hbar F_1^{(1)}(\rho(e_\mu)\eta)$ by computation.
\item $K=P^\dagger$ is Remark \ref{rem:c12-KP}, a cross-check; $K_1$ is computed by \eqref{eq:c12-master} at a finite endpoint.
\item The twist is Proposition \ref{prop:c12-twist} and Remark \ref{rem:c12-twistcmp}, a cross-check.
\end{itemize}
\end{remark}

\begin{remark}[Method of computation]
\label{rem:c12-method}
Every number is computed by two independent evaluations. The first evaluates each defect cycle in the scaled Schwinger form: with $t_e=\varepsilon\tau_e$, $\tau_\sigma\in[0,1]$ for the special edge and $\tau_e\in[1,\infty)$ for the other edges (both segments of a composite line included), the coefficient of $\varepsilon^0$ is a finite sum of the integrals
\[
J(n_t;\beta;N)=\int_0^1d\tau_\sigma\int_{[1,\infty)^{n_t-1}}\prod_{e\neq\sigma}d\tau_e\;\tau^\beta\,\Big(\sum_e\tau_e\Big)^{-N},
\]
$\tau^\beta$ a monomial, convergent for $N>|\beta|+n_t-1$, evaluated exactly in $\mathbb Q+\mathbb Q\log2+\mathbb Q\log3$ for $n_t\le3$ (Proposition \ref{prop:c12-closed}); letter derivatives through the generating functions of Proposition \ref{prop:c11-genfun}, the position through the injected momentum $r$ by exact interpolation on integer $r$, the fibre algebra in an explicit realisation of $V\oplus W_{\mathbb R}^\perp$ and $\Sigma$ over $\mathbb Q(i)$, colour in the realisations $\mathfrak{su}(2)$, $\mathfrak{su}(3)$ and $\mathfrak{gl}(3)$ by Lemma \ref{lem:c10-colour}, and exact linear algebra on bases fixed by representation counts. The second computes the renormalised amplitudes $W^{(1)}_{\rm ren}(x)$, $W^{(2)}_{\rm ren}(x,y)$ and the defect cycles in momentum space with an independently written two-component realisation of $S^{(1)}(x)$ (Definition \ref{def:c8-twocomp}), Feynman parameters for the edges $P$, the special edge as the factor $e^{-\varepsilon|p|^2}-1$, exact integrals where they close and otherwise thirty-digit numerics at several $\varepsilon$, extrapolated to $\varepsilon^0$ with the poles fitted, and along $K_b$ the endpoints $T=1,4,16$; it applies $Q$ to $W^{(1)}_{\rm ren}(x)$ (Theorem \ref{thm:c12-second}) and evaluates the second route of Theorem \ref{thm:c12-conj}. The two evaluations agree on every entry of Theorems \ref{thm:c12-ob1}, \ref{thm:c12-sol}, \ref{thm:c12-ar2}, \ref{thm:c12-F0} and of the first four rows of Table \ref{tab:c12-spectra}, and check the Wess--Zumino identities \eqref{eq:c12-wz} on every computed entry, the pole identities \eqref{eq:c12-pole}, \eqref{eq:c12-local2} and the position identities \eqref{eq:c12-pos}; the single-pair tables cover the $45$ pairs of slots with defect cycles with at most three letter derivatives in the first evaluation and one in the second. The overall sign of the defect cycles is fixed in each evaluation by a regression to the values of Chapters 4 and 9.
\end{remark}

\begin{proposition}[Closed forms]
\label{prop:c12-closed}
With $F_m$ the $m$-th simplex integrand of \eqref{eq:c12-special}:
\begin{itemize}
\item (i) for $n_t=2$, $16\pi^2a_0=\int_0^{1/2}F_0(\alpha)\log\frac{1-\alpha}\alpha\,d\alpha$ with $\alpha=\alpha_\sigma$, which is $F_0\log2$ for constant $F_0$;
\item (ii) for constant $F_0$ ($D=0$) and any $n_t$, $16\pi^2a_0=\frac{F_0}{(n_t-1)!}\log\frac{n_t}{n_t-1}$, so $\tfrac12F_0\log\tfrac32$ for $n_t=3$; for constant $F_1$ ($D=2$), $16\pi^2a_1=\frac{F_1}{n_t^2(n_t-1)(n_t-1)!}$;
\item (iii) for one vertex of $I$ and a composite line, $16\pi^2a_0=\int_0^{1/3}F_0(\beta)\big[(1-\beta)\log\frac{1-\beta}{2\beta}-(1-3\beta)\big]d\beta$, in $\mathbb Q+\mathbb Q\log2+\mathbb Q\log3$;
\item (iv) triangles of cubic vertices with the special edge have values in $\mathbb Q+\mathbb Q\log\tfrac32$.
\end{itemize}
\end{proposition}

\begin{proof}
(i) For $n_t=2$ the region where $\alpha_\sigma$ is the smaller coordinate is $\alpha\le\tfrac12$, with $\min_{f\neq\sigma}\alpha_f=1-\alpha$, and $\int_0^{1/2}\log\frac{1-\alpha}\alpha d\alpha=\log2$. (ii) By symmetry the region where $\alpha_\sigma$ is smallest is $1/n_t$ of the simplex, of volume $1/(n_t-1)!$, and the logarithmic average over it is $\log\frac{n_t}{n_t-1}$ (Theorem \ref{thm:c11-oddsimplex} (iv)). (iii) Proposition \ref{prop:c11-k1}. (iv) Theorem \ref{thm:c9-sing}.
\end{proof}

\begin{remark}[The computations that complete the construction]
\label{rem:c12-next}
The construction is complete through the two-leg parts of arity one along every generator and of arity two along $(\Delta,P_b)$, $(\Delta,\xi_\pm)$, $(\xi_+,\xi_-)$, $(\xi_\pm,\xi'_\pm)$, $(\xi_\pm,P_b)$, $(\xi_+,\eta_\pm)$, and through the components on the scalar bilinears and the spectra of the families. The next computations, each a finite evaluation of the classes of graphs above, are:
\begin{itemize}
\item on the action side, the two-leg arity-two equations along $(K_b,P_{b'})$, $(\eta,\eta')$ and $(\xi_-,\eta_\pm)$, which fix $t_K$ and give the one-loop forms of $[K,P]=D+M$ and $\{S,S\}=K$; the three- and four-leg parts of the arity-one equations, which fix the differences of the sector coefficients through Theorem \ref{thm:c12-cubic}, the four directions $G^K_j$, the free classes in the $\mathbf{15}$ of $\mathfrak{su}(4)$ at the weight of $(\xi_+,\xi_-)$, the arity-zero couplings of Proposition \ref{prop:c12-couplings} and the five-leg sector along $P_b$, and test the classical symmetry condition of Definition \ref{def:c12-cs}; the arity-three component of $s$ written term by term;
\item on the operator side, the identities of Lemma \ref{lem:c12-Cid} and Proposition \ref{prop:c12-typeI} as computed identities on the composite loops with two and three legs and on the scalar bilinears, \eqref{eq:c11-logcov} and the full pole identity on the line rows; $Q_1$, $S_1$, $M_1$, $K_1$ and $P_1$ along generic $b$ on the families at generic kinematics; the arity-two operator kernels of Proposition \ref{prop:c12-master2} and the two-sided relations (R2)--(R6); the second evaluation on the families with two gaugino components, the gluon sector and the mixed family; $\mathrm{Gram}_1$ from the two-point graphs for Remark \ref{rem:c12-KP}; the overall sign of the defect cycles from the sign rule of Definition \ref{def:c11-defect} (ii);
\item for the twist, the transport map $\Phi$ with the normalisation of [BK] and the $(\xi,\xi)$ kernels of the entries $b\,c$ and $b\,b$ (Remark \ref{rem:c12-twistcmp}).
\end{itemize}
With the arity-two kernels and the components on the weight-$\tfrac52$ descendants, (R3) on the trace bilinear becomes the Konishi relation between $Q_1$, $S_1$ and $D_1=-48$ units, a scheme-independent test of the combinations \eqref{eq:c12-F0inv}.
\end{remark}

\begin{remark}[Two-component forms]
\label{rem:c12-twocomp}
With Definition \ref{def:c8-twocomp}, $\sigma^m_{\alpha\dot\alpha}=(\vec\sigma,i)$, $\bar\sigma^{m\dot\alpha\alpha}=(-\vec\sigma,i)$, $\varepsilon^{12}=\varepsilon^{\dot1\dot2}=+1$, $A_m=-igA^a_mT^a$, $A_I=-igA^a_IT^a$, $A_{ij}=A_I\Sigma^I_{ij}$, $A^{ij}=\tfrac12\varepsilon^{ijkl}A_{kl}$, $\lambda=-\tfrac{ig}{\sqrt2}\lambda^aT^a$, $\operatorname{tr}(T^aT^b)=\tfrac12\delta^{ab}$, the unit $-\tfrac{g^2N}{32\pi^2}$, and $(-,-)$ the pairing of $\mathfrak g$ around the components:
\begin{itemize}
\item Pairings. $\int(A,A^*)=\int d^4x\{(A_m,A^*_m)+\tfrac14(A_{ij},A^{*ij})\}$, $\int(\lambda^*,\lambda)=\int d^4x\{(\lambda_\alpha{}^i,\lambda^{*\alpha}{}_i)-(\bar\lambda_{\dot\alpha i},\bar\lambda^{*\dot\alpha i})\}$, so $\mathcal N=\int d^4x\{2(A_m,A^*_m)+\tfrac12(A_{ij},A^{*ij})-4(\lambda_\alpha{}^i,\lambda^{*\alpha}{}_i)+4(\bar\lambda_{\dot\alpha i},\bar\lambda^{*\dot\alpha i})-(c^*,c)\}$, with the order of the odd factors of Lemma \ref{lem:c6-count}.
\item Supercharge data. For $\xi=(\xi_\alpha{}^i,0)\in S_+\otimes W$: $\Gamma_V^m=i\,\bar\lambda_{\dot\alpha i}\bar\sigma^{m\dot\alpha\alpha}\xi_\alpha{}^i$, $(\Gamma_\perp)_{kl}=2\varepsilon_{jikl}\,\xi_\alpha{}^j\varepsilon^{\alpha\beta}\lambda_\beta{}^i$, $(\xi,\lambda^*)=\xi_\alpha{}^i\lambda^{*\alpha}{}_i$, and
\[
\begin{aligned}
&T_0=-\int d^4x\,\big(\partial_mA_m,\xi_\alpha{}^i\lambda^{*\alpha}{}_i\big),\quad T_1=-i\int d^4x\,\big(\bar\lambda_{\dot\alpha i}\bar\sigma^{m\dot\alpha\alpha}\xi_\alpha{}^i,A^*_m\big),\quad T_2=-\tfrac12\int d^4x\,\varepsilon_{jikl}\big(\xi_\alpha{}^j\varepsilon^{\alpha\beta}\lambda_\beta{}^i,A^{*kl}\big),\\
&T_3=-\tfrac14\int d^4x\,\big((\partial_mA_n-\partial_nA_m)\sigma^m_{\alpha\dot\alpha}\bar\sigma^{n\dot\alpha\beta}\xi_\beta{}^i,\lambda^{*\alpha}{}_i\big),\quad T_4=-\tfrac i2\int d^4x\,\big(\varepsilon^{\alpha\beta}\xi_\beta{}^j\sigma^m_{\alpha\dot\alpha}\partial_mA_{ij},\bar\lambda^{*\dot\alpha i}\big),
\end{aligned}
\]
$g_1=\int d^4x\,(\Gamma_V^n,-\partial^2A_n+\partial_n\partial_mA_m)$, $g_2=-\int d^4x\,(\Gamma_V^n,\partial_n\partial_mA_m)$, $g_3=-\tfrac14\int d^4x\,((\Gamma_\perp)_{kl},\partial^2A^{kl})$, $g_4=-\int d^4x\,(\partial^2c,\xi_\alpha{}^i\lambda^{*\alpha}{}_i)$. The terms with $\bar\xi_{\dot\alpha i}$ follow by $\lambda_\alpha{}^i\leftrightarrow\bar\lambda_{\dot\alpha i}$, $\lambda^{*\alpha}{}_i\leftrightarrow\bar\lambda^{*\dot\alpha i}$, $\sigma\leftrightarrow\bar\sigma$.
\item Special supercharge data. For $\eta=(\eta^\alpha{}_i,0)$: $(\rho(e_m)\eta)_{\dot\alpha i}=i\,\eta^\alpha{}_i\sigma^m_{\alpha\dot\alpha}$, and for $\eta=(0,\bar\eta^{\dot\alpha i})$: $(\rho(e_m)\eta)_\alpha{}^i=-i\,\sigma^m_{\alpha\dot\alpha}\bar\eta^{\dot\alpha i}$ \eqref{eq:c8-pairing}; $s_0=\int d^4x\,(A^{kj}\varepsilon_{\alpha\beta}\eta^\beta{}_j,\lambda^{*\alpha}{}_k)$ and $\int(\rho(A|_V)\eta,\lambda^*)=-i\int d^4x\,(A_m\eta^\alpha{}_i\sigma^m_{\alpha\dot\alpha},\bar\lambda^{*\dot\alpha i})$.
\item Sector parts. $S^{(1)}(\Delta)_\perp=\tfrac14\int d^4x\,((2x_n\partial_n+2)A_{ij},A^{*ij})$, $S^{(1)}(\Delta)_\lambda=-\int d^4x\{((2x_n\partial_n+3)\lambda_\alpha{}^i,\lambda^{*\alpha}{}_i)-((2x_n\partial_n+3)\bar\lambda_{\dot\alpha i},\bar\lambda^{*\dot\alpha i})\}$, $S^{(1)}(\Delta)_V=\int d^4x\{((2x_n\partial_n+2)A_m,A^*_m)-(2x_n\partial_nc,c^*)\}$, and $S^{(1)}(P_b)_d$ the same with $2x_n\partial_n+2w$ replaced by $b_n\partial_n$; $S^{(1)}(M_\omega)_{\rm spin}=\int d^4x\{\omega_{mn}(A_m,A^*_n)+\tfrac14\omega_{mn}[(\lambda^{*\alpha}{}_i(\sigma^m\bar\sigma^n)_\alpha{}^\beta\lambda_\beta{}^i)-(\text{the same on the dotted components with }\bar\sigma^m\sigma^n)]\}$.
\item Obstructions. $\hbar\,\mathrm{Ob}^{(0)}|_2=-\tfrac{27}8\int d^4x\,(\partial^2c,\partial_mA_m)$ units, $\hbar F_1^{(0)}|_2=\tfrac{27}{16}\int d^4x\,(\partial_mA_m,\partial_nA_n)$ units, and
\[
\hbar\,\mathrm{Ob}^{(1)}(\Delta)|_2=\int d^4x\Big\{-\tfrac{11}3\Big[\tfrac12(F^{\rm lin}_{mn},F^{\rm lin}_{mn})+\tfrac14(\partial_mA^{ij},\partial_mA_{ij})\Big]+\tfrac{94}3\,i\big(\bar\lambda_{\dot\alpha i},\bar\sigma^{m\dot\alpha\alpha}\partial_m\lambda_\alpha{}^i\big)-\tfrac{14}3(\partial_mc,A^*_m)\Big\}\ \text{units},
\]
$F^{\rm lin}_{mn}=\partial_mA_n-\partial_nA_m$; along $\xi$ the obstruction is \eqref{eq:c12-ob1} with $g_1,\dots,g_4$ above.
\item Solutions. $\hbar F_1^{(1)}(\Delta)|_2=\int d^4x\{\tfrac{22}3[(A_m,A^*_m)+\tfrac14(A_{ij},A^{*ij})]-\tfrac{47}3[(\lambda_\alpha{}^i,\lambda^{*\alpha}{}_i)-(\bar\lambda_{\dot\alpha i},\bar\lambda^{*\dot\alpha i})]-\tfrac83(c^*,c)\}$ units plus the free directions; $\hbar F_1^{(1)}(\xi)|_2=(\log\tfrac32-\tfrac{79}{12})(T_1+T_2)$ units $+\,t_\xi(T_1+T_2+T_3+T_4)$ units; along $K_b$, $\int(\iota_bA)(\iota_xA^*)=\int d^4x\,(b_mA_m,x_nA^*_n)$ and $\rho_\mu\rho_\nu$ is $\sigma^\mu\bar\sigma^\nu$ on the undotted and $\bar\sigma^\mu\sigma^\nu$ on the dotted components.
\item Arity two. $S^{(2)}(\xi,\xi')=\int\{\tfrac14(\Gamma(\xi,\xi'),\Gamma(\lambda^*,\lambda^*))-\tfrac12(\xi,\lambda^*)(\xi',\lambda^*)-(\iota_{\Gamma(\xi,\xi')}A,c^*)\}$ with $\Gamma(\xi,\bar\xi')|_V^m=i\,\bar\xi'_{\dot\alpha i}\bar\sigma^{m\dot\alpha\alpha}\xi_\alpha{}^i$ for $\xi\in S_+\otimes W$, $\bar\xi'\in S_-\otimes W^*$, and $(\Gamma(\xi,\xi')|_\perp)_{kl}=2\varepsilon_{jikl}\,\xi_\alpha{}^j\varepsilon^{\alpha\beta}\xi'_\beta{}^i$ for $\xi,\xi'\in S_+\otimes W$; $\psi_1=\int d^4x\,x_m(\Gamma^m_V,c^*)$, $\psi_2=\int d^4x\,x_m(A^*_m,\xi_\alpha{}^i\lambda^{*\alpha}{}_i)$.
\item Bilinears. For $O=\tfrac1{16}\int\omega\,w_1^{ij}w_2^{kl}(A_{ij},A_{kl})$: $\mathcal B_1=\tfrac14\int\omega\,[(w_1^{kl}A_{kl},w_{2\,ji}\xi^{\alpha i}\lambda_\alpha{}^j)+(1\leftrightarrow2)]$, $\mathcal B_2=\tfrac14w_1^{ij}w_{2\,ij}\int\omega\,(A_{ji},\xi^{\alpha i}\lambda_\alpha{}^j)$, $\mathcal B_3=\tfrac14w_1^{ij}w_{2\,ij}\int\omega\,i(\bar\lambda_{\dot\alpha i},A_m\bar\sigma^{m\dot\alpha\alpha}\xi_\alpha{}^i)$, $\mathcal B_4=\tfrac14w_1^{ij}w_{2\,ij}\int\omega\,(c,\xi_\alpha{}^i\lambda^{*\alpha}{}_i)$; Theorem \ref{thm:c12-F0} then reads $Q_1(O)=(\tfrac{79}{12}-3\log\tfrac32-t_\xi)\mathcal B_1+4\log\tfrac32(\mathcal B_2+\mathcal B_3)$ in units $-\tfrac{g^2N}{32\pi^2}$, and on the traceless bilinears the transformation along $Q^i_\alpha$ is $1+c_Q$ times the classical one.
\item Generators and directions. $B$ is $+1$ on $\xi_\alpha{}^i$, $\bar\eta^{\dot\alpha i}$ and $-1$ on $\bar\xi_{\dot\alpha i}$, $\eta^\alpha{}_i$ (Definition \ref{def:c8-B}); the residual freedom $t=aB+\tfrac c2\lambda$ multiplies the one-loop actions of $Q^i_\alpha$, $\bar Q_{\dot\alpha i}$, $S^\alpha{}_i$, $\bar S^{\dot\alpha i}$, $P_m$, $K_m$ by $1-(a-\tfrac c2)$, $1-(-a-\tfrac c2)$, $1-(-a+\tfrac c2)$, $1-(a+\tfrac c2)$, $1+c$, $1-c$ at one-loop order; $Q_1(O)=\xi_\alpha{}^iQ_1(O)^\alpha{}_i+\bar\xi_{\dot\alpha i}\bar Q_1(O)^{\dot\alpha i}$ and the same with $\eta^\alpha{}_i$, $\bar\eta^{\dot\alpha i}$ along $S^\alpha{}_i$, $\bar S^{\dot\alpha i}$.
\item Position decompositions. The position part of $S_1$ along $S^\alpha{}_i$ at $x_O$ is $Q_1$ along $\bar\xi_{\dot\alpha i}=i\,x_{O,m}\eta^\alpha{}_i\sigma^m_{\alpha\dot\alpha}$, and along $\bar S^{\dot\alpha i}$ it is $Q_1$ along $\xi_\alpha{}^i=-i\,x_{O,m}\sigma^m_{\alpha\dot\alpha}\bar\eta^{\dot\alpha i}$; $\Upsilon_{K_b}=v^m(x_O)\Upsilon_{P_m}-(b\cdot x_O)\Upsilon^0_\Delta+\Upsilon^0_{M_{\omega_{x_O}}}+\Upsilon^0_{K_b}$ with $v^m(x_O)=x_O^2b^m-2(b\cdot x_O)x_O^m$ and $(\omega_{x_O})^m{}_n=2(b^mx_{O,n}-x_O^mb_n)$.
\item Weinberg. On the Dirac column $(\lambda_\alpha{}^i,\bar\lambda^{\dot\alpha}{}_i)$, $\rho(v)=\gamma_W^\mu v_\mu$ gives Weinberg's form of each line (Proposition \ref{prop:c8-twocomp}).
\end{itemize}
\end{remark}
